%%
%% fall-lecture18.tex
%% 
%% Made by Alex Nelson <pqnelson@gmail.com>
%% Login   <alex@lisp>
%% 
%% Started on  2026-11-08T10:40:14-0800
%% Last update 2026-11-08T10:40:14-0800
%% 

\lecture[Measure spaces]{}

\begin{definition}
A \define{Measurable Space} is a pair $(X,\mathcal{A})$ where
$\mathcal{A}$ is a $\sigma$-algebra on $X$.
\end{definition}

\begin{definition}
Let $X$ be a set.
\begin{enumerate}
\item An \define{Algebra (of sets)} of $X$ is a collection
  $\mathcal{A}\subset\powerset{X}$ of subsets of $X$ which is closed
  under finite unions and complements.
\item A \define{$\sigma$-Algebra} of $X$ is an algebra $\mathcal{A}$
  of $X$ which is closed under \emph{countable} unions and complements.
\end{enumerate}
\end{definition}

\begin{xca}
Let $\mathcal{A}$ be an algebra of $X$.
Prove or find a counter-example: if for each countable family
$A_{k}\in\mathcal{A}$ of disjoint subsets (so $A_{k}\cap A_{j}=\emptyset$ for all $j\neq k$)
we have their union belong to $\mathcal{A}\ni\bigcup_{k}A_{k}$, then
$\mathcal{A}$ is a $\sigma$-algebra.
\end{xca}

\begin{notation}
We will commonly abuse notation and write $\bigsqcup_{\lambda}A_{\lambda}$
to indicate $\bigcup_{\lambda}A_{\lambda}$ under the assumption that
$A_{\lambda}\cap A_{\lambda'}=\emptyset$ for all $\lambda\neq\lambda'$.
\end{notation}

\begin{example}
Let $X=[0,1)$ be the half-open interval.
Let $\mathcal{A}=\{[0,1/2),[1/2,1),\emptyset,X\}$. Then $\mathcal{A}$
    is an algebra.
We can consider more generally the dyadic intervals which are
subintervals of $X$:
\begin{equation}
\mathcal{A}_{n}=\{k+[0,2^{-n})\mid 0\leq k\leq 2^{n}\}\cup\{\emptyset,X\}
+\mbox{(their unions)}.
\end{equation}
These are a monotone increasing family of algebras
$\mathcal{A}_{1}\subset\mathcal{A}_{2}\subset\dots$. These are all
$\sigma$-algebras. This describes (in probability theory) tossing a
coin infinitely many times, $\mathcal{A}_{n}$ describes the outcomes
of the first $n$ tosses.
\end{example}

\begin{example}
A continuous version of the previous example would be $C[0,\infty)=X$
or if we do not require continuity (of elements of $X$) we could take $X=\RR^{[0,\infty)}$.
This describes elements $x_{t}$ or $x(t)$. Then $\mathcal{A}_{t}$ is
the increasing family of events---events determined by the motion up
to time $t$. Think of Brownian motion.
\end{example}

\begin{lemma}
Intersections of $\sigma$-algebras of $X$ is a $\sigma$-algebra of $X$.
\end{lemma}

\begin{lemma}
If $\mathcal{F}$ is a family of subsets of $X$ (so
$\mathcal{F}\subset\powerset{X}$), then we can consider the smallest
$\sigma$-algebra of $X$ containing $\mathcal{F}$. This is denoted
$\sigma(\mathcal{F})$ (it is also considered the $\sigma$-algebra
``generated'' by the elements of $\mathcal{F}$).
\end{lemma}

\begin{definition}
We call $\sigma(\mathcal{F})$ the \define{$\sigma$-Algebra Generated by $\mathcal{F}$}.
\end{definition}

\begin{definition}
Let $(X,\topology)$ be a topological space. Then the \define{Borel $\sigma$-Algebra}
of $(X,\topology)$, denoted either $\mathcal{B}_{X}$ or something similar, is the $\sigma$-algebra generated by $\topology$:
$\mathcal{B}_{X}=\sigma(\topology)$.
\end{definition}

\begin{definition}
Elements of $\mathcal{B}_{\RR}$ (the Borel $\sigma$-algebra for $\RR$
equipped with standard topology) are called \define{Borel Sets}.
\end{definition}

\begin{example}
\begin{enumerate}
\item Open subsets of $\RR$ are Borel sets
\item Closed subsets of $\RR$ are Borel sets
\item $G_{\delta}$ subsets of $\RR$ are Borel sets
\item $F_{\sigma}$ subsets of $\RR$ are Borel sets
\end{enumerate}
\end{example}

\begin{xca}
Show the Borel algebra $\mathcal{B}_{\RR}$ is generated by half-open
intervals $[a,b)$ and closed intervals $[a,b]$. Show
  $\mathcal{B}_{\RR}$ is generated by open (resp., closed) rays
  $(a,\infty)$ and $(-\infty,b)$ (resp.,  $[a,\infty)$ and $(-\infty,b]$).
\end{xca}

\begin{definition}
Let $(X,\mathcal{A}_{X})$ and $(Y,\mathcal{A}_{Y})$ be sets equipped
with $\sigma$-algebras. We can define the \define{Product $\sigma$-Algebra}
denoted $\mathcal{A}_{X}\otimes\mathcal{A}_{Y}$ as the
$\sigma$-algebra on $X\times Y$ generated by elements $E_{X}\times E_{Y}\in\mathcal{A}_{X}\times\mathcal{A}_{Y}$.
In other words, $\mathcal{A}_{X}\otimes\mathcal{A}_{Y}:=\sigma(\mathcal{A}_{X}\times\mathcal{A}_{Y})$.
\end{definition}

\begin{node}
This means that $\mathcal{B}_{\RR^{2}}$ is generated by rectangles in $\RR^{2}$.
\end{node}

\begin{node}
Of course, we can generalize product $\sigma$-algebras to finite products.
\end{node}

\begin{xca}
For infinite products $(X_{\alpha},\mathcal{A}_{\alpha})_{\alpha\in A}$ we have
$\bigoplus_{\alpha\in A}\mathcal{A}_{\alpha}$ which is a
$\sigma$-algebra on $\prod_{\alpha\in A}X_{\alpha}$. Try defining what
this should be.
\end{xca}

\subsection{Measures on $(X,\mathcal{A})$}

\begin{definition}
A \define{Measure} is a map $\mu\colon\mathcal{A}\to[0,\infty]$ such
that
\begin{enumerate}
\item $\mu(\emptyset)=0$
\item\textsc{$\sigma$-additivity} (or \textsc{Countable additivity}): If we have a [countable family of] disjoint
  $\sigma$-measurable sets $E_{j}\in\mathcal{A}$, then
  \begin{equation}
\mu\left(\bigsqcup_{j}E_{j}\right)=\sum_{j}\mu(E_{j}).
  \end{equation}
  (This depends critically on \emph{countably} many $E_{j}$.)
\end{enumerate}
\end{definition}

\begin{node}
We should think of measures as mass, volume, or length (depending on
the situation).
\end{node}

\begin{definition}
A \define{Measure Space} is a triple $(X,\mathcal{A},\mu)$ where $X$
is a set, $\mathcal{A}$ is a $\sigma$-algebra of $X$, and
$\mu\colon\mathcal{A}\to[0,\infty]$ is a measure.
\end{definition}

Now, let us look at some properties of measures.

\begin{proposition}
Measures are monotone, subadditive, and ``continuous'' [in some
  appropriate sense].

Monotone: if $E\subset E'$ are measurable sets, then $\sigma(E)\leq\sigma(E')$.

Subadditivity: $\mu(\bigcup_{j}E_{j})\leq\sum_{j}\mu(E_{j})$ where
$E_{j}$ are not necessarily disjoint.

Continuous: If $E_{n}\uparrow E$ is an increasing [convergent] sequence of
measurable sets, then $\mu(E_{n})\to\mu(E)$.

Continuous 2: If $E_{n}\downarrow E$ is a decreasing [convergent]
sequence of measurable sets \emph{and also} $\mu(E_{1})<\infty$ is
finite measure, then $\mu(E_{n})\to\mu(E)$.
\end{proposition}

\begin{proof}
Let us prove continuity. Let $E_{n}\uparrow E$ be an increasing
monotone sequence of measurable sets. (If $E_{0}$ is not defined, set $E_{0}=\emptyset$.) Then
\begin{subequations}
  \begin{align}
\mu\left(\bigcup_{n}E_{n}\right) &= \mu\left(\bigsqcup_{n=1}(E_{n}\setminus E_{n-1})\right)\\
&= \sum_{n=1}\bigl(\mu(E_{n})-\mu(E_{n-1})\bigr)\\
&=\lim_{N\to\infty}\sum^{N}_{n=1}\bigl(\mu(E_{n})-\mu(E_{n-1})\bigr)\\
&=\lim_{N\to\infty}\mu(E_{N})\quad\mbox{by telescoping sum}.
  \end{align}
\end{subequations}
Hence meriting the name ``continuity''.
\end{proof}

\begin{puzzle}
In proof of ``continuity 2'', when is $\mu(E_{1})<\infty$ needed?
\end{puzzle}

\begin{example}[Delta measures]
Let $x\in X$. The function $\delta_{x}\colon\powerset{X}\to[0,\infty]$
defined by
\begin{equation}
\delta_{x}(E)=\begin{cases}1 & x\in E\\
0 & \mbox{otherwise}
\end{cases}
\end{equation}
is a measure called the \define{$\delta$-Measure}.
\end{example}

\begin{example}[Counting measure]
We could define $\mu(E)$ equal to the number of elements in $E$ (or
$\mu(E)=\infty$ if $E$ is an infinite set). This gives us a measure on
$(X,\powerset{X})$. 
\end{example}

\begin{node}
Further examples are not so easy, so let us state the fundamental theorem.
\end{node}

\begin{theorem}[Carath\'{e}odory]
Let $\mathcal{R}$ be an algebra and let
$\lambda\colon\mathcal{R}\to[0,\infty]$ be such that
\begin{enumerate}
\item (``$\sigma$-additivity
  when it makes sense) Suppose whenever $E=\bigsqcup E_{j}$ is a countable union of disjoint
  $E_{j}\in\mathcal{R}$, and when $E\in\mathcal{R}$, then we have
  $\lambda(E)=\sum_{j}\lambda(E_{j})$. In this situation, the
  Carath\'{e}odory theorem asserts $\lambda$ extends to a measure $\mu$.
\item Moreover, this extension is unique if $\lambda$ is $\sigma$-finite.
\end{enumerate}
Such $\lambda$ are called \define{Premeasures} (or, confusingly,
\define{Measures on Algebras}).
\end{theorem}